\paragraph{Protocol fragility.} The pilot study trained one deep hedger (arm A0) on a single
pre-2016 split with one seed. It reported a CVaR reduction of $-0.25$ at \FocusBps{} bps, significant
under the month-clustered bootstrap. Under walk-forward retraining the same arm gives \DiffAZeroFocus{}
(95\% CI \CIAZeroFocus{}) with \NSeeds{} seeds, and only $-0.05$ with the first three seeds.
Its seed-to-seed standard deviation is 0.05--0.15 across arms, the same order as the effect.
After the Romano--Wolf correction over all \NFamily{} comparisons its adjusted $p$-value is \PRWAZeroFocus{}.
A single split with a single seed therefore overstates both the size and the reliability of the gain.

\paragraph{Realistic costs: a three-parameter correction is enough.} At \FocusBps{} bps the linear
correction L reduces CVaR by \DiffLFocus{} (CI \CILFocus{}, Romano--Wolf $p=$\PRWLFocus{}). It beats delta
in 10 of 11 test years and is the only method significant after correction at 10, 15 and 25 bps.
The best deep arms (A0 and A5) reach the same point estimate, but with wider intervals and no
significance after correction (Table~2). In money terms L saves about HK\$1{,}400 of tail loss per HK\$1m of
option notional, 2.3\% of the average premium (Table~3).

\paragraph{High costs: deep hedging pays.} At 50 bps the deep arms trained on GBM
(\DiffAZeroFifty{} for A0, \DiffAOneFifty{} for A1) and A5 (\DiffAFiveFifty{}) cut tail risk by 0.30--0.37,
about 5--6\% of the premium, roughly twice as much as L (\DiffLFifty{}). A0 remains significant after the
Romano--Wolf correction ($p=$\PRWAZeroFifty{}); A1 and A5 are borderline (\PRWAOneFifty{}, \PRWAFiveFifty{}).
This matches reports that deep-hedging gains grow with transaction costs \citep{francois2025ivs,bracha2025drl}.

\paragraph{Training data: realism does not help.} The block bootstrap reproduces the real data's tails
and volatility clustering closely (excess kurtosis 7.7 vs 6.7; Hill index 3.4 vs 3.5), while GBM has
none. Yet training on it (A3: \DiffAThreeFocus{}) or on real windows (A4: \DiffAFourFocus{}) is no better
than training on GBM. GARCH-$t$ fitted to pre-year data overshoots the tails (kurtosis 136) and makes
training unstable; capping the tails (A2b, post hoc) stabilises it but does not beat delta. This agrees
with \citet{hashimoto2026realism}: fidelity of the generator is not what drives hedging performance.

\paragraph{Hull--White form.} Fitted to the variance of hedging P\&L, the minimum-variance form
of \citet{hull2017optimal} \emph{increases} CVaR relative to delta at every cost. Its $n(d_1)$ scaling
puts the correction where the option's gamma is largest. Without traded option prices, there is no
spot--implied-volatility co-movement for it to capture, so the variance-optimal fit does not carry over
to the tail. We read this as a limitation of synthetic option prices rather than evidence against Hull--White.
